\section{Discussion}
\label{sec:discussion}

\subsection{Why binding helps, and why it is not enough}

Physical binding helps because it changes the coordinate system of the state that attackers observe. Plaintext cache attacks depend on alignment: inversion assumes that cached values live in the public projection basis, collision assumes that locally generated rows can be compared to leaked rows, and injection assumes that a standard model can interpret the stolen K/V history. Orthogonal \sys{} breaks these assumptions by storing K/V in a session-specific basis while preserving attention for the device that generated the basis.

The same experiments show why binding alone is not enough, and exactly which property a software-key obfuscator lacks that physical binding supplies. Our migration test makes the distinction concrete: a software obfuscator protects the cache only while its key stays secret, so an adversary who dumps the cache and the key from the same process recovers everything; \sys{} has no such key, so a copied cache plus the entire software stack stays bound to the originating device. This is the core reason to use a PUF, and it is a property of where the secret lives, not of avoiding decryption --- indeed our strongest confidentiality variant, affine masking, deliberately decrypts inside attention.

Binding is also not the same as confidentiality. An orthogonal map preserves not just row norms but every right-invariant statistic of the cache: the norms are only the diagonal of the Gram matrix $XX^\top$, which is itself preserved, as is the singular-value spectrum. Finite-candidate attackers can compare any of these without learning the session basis, and all three recover protected secrets at top-1 1.000. The design problem is therefore sharper than ``make the cache device-bound'': a protected cache must remove the whole basis-invariant class, which is what motivates the additive, non-orthogonal affine mask.

\subsection{Why fixed bases fail}

A natural conjecture is that a PUF-derived basis is secure because the PUF root is physically unclonable. Our profiling results reject that conjecture for fixed sessions. If an adversary obtains aligned plaintext/protected pairs from the same session, the adversary does not need to clone the PUF; the adversary can learn the induced orthogonal map directly. This distinction reframes the design problem: the PUF root provides device binding, but session refresh provides profiling resistance.

The invariant-class result rejects a second conjecture: that any unknown orthogonal basis is sufficient for cache-only confidentiality. Norm, Gram-matrix, and singular-spectrum matchers each recover protected finite secrets with top-1 accuracy 1.000 on Qwen3-0.6B and Llama-3.2-1B. This is why we treat orthogonal \sys{} as a binding baseline and affine masking as the current confidentiality-oriented redesign.

\subsection{Why affine masking is the current redesign candidate}

Unit-normalizing exported cache rows blocks the norm attack in our prototype, but it requires a trusted sidecar for the original norms. That sidecar weakens the system story: if the deployment already has trusted storage for per-row metadata, reviewers can ask why the KV cache itself is not stored there. Affine masking avoids that objection. The mask is not stored as content metadata; the device regenerates it from the PUF root, session nonce, layer, head, and position.

Affine masking is not free. It changes the orthogonal-only property that attention can operate directly on protected coordinates; the legitimate path now regenerates and subtracts masks inside attention. It also introduces cancellation risk, visible in the Qwen3 fp32 sanity gap increasing from roughly $3\times10^{-5}$ for the orthogonal path to $4.90\times10^{-4}$ at mask std128, and it roughly doubles decode overhead. The mask strength is a controllable tradeoff: sweeping it shows the strongest invariant (Gram) attack falling from $0.975$ at std~4 to near the $1/32$ chance level by std~64--128, while the fp32 error grows monotonically, so std~128 is the knee rather than a tuned constant. The security gain is correspondingly larger: it removes the whole invariant class (Qwen3 Gram/spectrum/norm top-1 to $0.033/0.033/0.050$). These numbers make affine masking the best current direction, not a completed production design.

\subsection{Precision is a systems requirement}

The fp32 results show exact equivalence, whereas bf16 long-decode experiments expose early divergence on two prompts. This result does not invalidate the algebraic construction; it shows that protected-coordinate arithmetic can redistribute channel magnitudes and thereby change low-precision rounding trajectories. For deployment, the relevant question is not whether the identity holds in real arithmetic, but whether the inference stack preserves enough precision for the identity to remain behaviorally stable. Affine masking raises this requirement further because large masks must be subtracted before attention. We therefore treat fp32 attention activations, fp32 cache storage, stochastic rounding, or masking families optimized for quantized inference as implementation choices rather than optional refinements.

Rather than treating cache protection as a storage-layer encryption problem, we argue that LLM serving systems should treat persistent attention state as session-bound computation state. This prescription changes which secrets must survive compromise: copied tensors should not be enough to reconstruct or replay the original context.

\subsection{Compatibility with fused attention kernels}

Because the cache-migration surface is created by optimized serving systems --- paged attention and KV offloading --- a defense must eventually run inside their fused kernels rather than in eager attention. The orthogonal path composes cleanly with the standard fused entry point: the rotations are linear maps applied to $q$, $k$, and $v$ \emph{around} attention, so routing the rotated tensors through \texttt{scaled\_dot\_product\_attention} reproduces the plain logits to within $4\times10^{-5}$, and the $O_k$/$O_v$ multiplications fold naturally into the QKV and output projections. In fp32 this dispatches to the math/memory-efficient backend; fp16 FlashAttention additionally requires the precision care discussed below. The affine-mask path is the harder case, because the mask must be subtracted between the cache read and the QK/AV products. We implemented a fused Triton flash-decoding kernel that regenerates the PUF mask in-register via the same counter-based Philox RNG as the write path and subtracts it in the SRAM epilogue of the KV load, so de-masking adds no HBM traffic and no separate pass (RQ12, Table~\ref{tab:fused-kernel}). A split-K design (grid $(B,H_q,n_{\text{splits}})$ plus a log-sum-exp reduction) lifts SM occupancy and eliminates the long-context regression a naive single-program-per-head kernel suffers. The fused output is token-identical to the eager de-mask path on both Qwen3-0.6B and Llama-3.2-1B and is $4$--$5$ percentage points cheaper in decode overhead ($\approx{-}29\%$ vs.\ plain, against $\approx{-}33\%$ for eager). The precision concern is also milder than the bf16 result suggests: fp16, the precision FlashAttention typically uses, diverges on only $14.1\%$ of continuations for the fused affine path versus $20.3\%$ for both orthogonal and eager-affine (all at 64 prompts and 64 tokens, fp32 cache; the shorter 32-token diagnostic in RQ3 gives $15.6\%$/$56.25\%$ for fp16/bf16 respectively---the rate grows with decode length but the ordering is stable), because fp16's larger mantissa resolves the rotated channel magnitudes and the fused kernel keeps the entire de-mask in fp32. Paged and flash-attention integration of the fused kernel remains engineering work.

\subsection{Scope: where physical binding is the right tool}

\sys{} is not a drop-in replacement for confidential-computing enclaves. Its niche is edge and on-premise inference where a trusted execution environment is unavailable or unattested but the device has a usable PUF/root-of-trust, and where the realistic threat is post-compromise migration of persistent cache state rather than a live attacker inside the running enclave. In that setting the alternatives are weaker: software-key obfuscation migrates with the cache, and full KV encryption requires a device-bound key whose lifecycle must itself be solved. Where a strong attested TEE is available, it subsumes much of this; \sys{} targets the large deployment surface where it is not.

\subsection{Limitations}

The current prototype has five primary limitations. First, orthogonal \sys{} is not a complete confidentiality mechanism because the full basis-invariant class --- norm, Gram, and spectrum matching --- recovers protected secrets with top-1 accuracy 1.000. Second, affine masking removes that class and is now utility-validated (fp32-exact on perplexity, HellaSwag, MMLU, and long decoding) with a fused split-K de-masking kernel validated at $\approx 29\%$ decode overhead (RQ12), but it remains fp32-only and the kernel targets contiguous-cache decode rather than paged/flash serving. Third, the square-projection inversion attack that gives the dramatic top-1 1.000 token-recovery baseline is specific to Qwen3-0.6B; on the other families we rely on injection and finite-candidate matching, which generalize but are less visually striking, so the cross-model leakage story rests on those rather than on inversion. Fourth, the PUF is simulated, so the current work validates the software interface and noise model but not physical reliability, uniqueness, or helper-data leakage. Fifth, the fused kernel is benchmarked on the decode path in isolation; full integration with paged-attention and flash-attention serving stacks (including prefill fusion and bf16/fp16 policies) remains engineering work.

These limitations define the next mechanism rather than invalidate the direction. Non-square projection attacks can use injection, finite-candidate matching, and task-level leakage metrics rather than square V-inversion. Physical PUF traces can replace the simulator root while preserving the HMAC context derivation. Serving benchmarks can measure the cost of fused rotation and mask-subtraction kernels under bf16, fp32, and fp32-cache policies.

\section{Related Work}
\label{sec:related}

\textbf{KV-cache privacy attacks and defenses.} Shadow in the Cache identifies inversion, collision, and injection attacks against plaintext KV-caches and proposes KV-Cloak as a software-secret obfuscation defense \cite{shadow_cache}. Our work builds on that threat surface but changes the trust boundary: \sys{} binds the cache representation to a physical session basis rather than to exportable software matrices. The broader class of microarchitectural cache side-channel attacks, such as \textsc{Flush+Reload}~\cite{flush_reload}, motivates treating leaked cache state as an exploitable observation surface rather than inert storage.

\textbf{Efficient LLM serving.} KV-cache systems, paged attention, and cache offloading improve throughput by externalizing persistent inference state. This performance optimization inadvertently creates the cache-at-rest and cache-migration surface studied in this paper. Orthogonal \sys{} preserves the reuse semantics of the cache instead of encrypting and decrypting it around every attention operation; affine masking preserves reuse but requires mask subtraction inside attention. Systems such as PagedAttention/vLLM~\cite{pagedattention} and offloading engines like FlexGen~\cite{flexgen} externalize and migrate this persistent state for throughput, which is precisely what enlarges the cache-migration attack surface.

\textbf{Hardware roots of trust and PUFs.} PUFs have been used for device authentication and key derivation. \sys{} uses the PUF root differently: the PUF determines the coordinate system and affine masks in which attention state is stored. This makes PUF security not only an identity property but also a state-interpretation property. Our construction builds on classic silicon PUFs~\cite{gassend_silicon} and PUF-based key generation~\cite{suh_devadas_puf}, together with fuzzy extractors~\cite{dodis_fuzzy} for stable root reconstruction from noisy responses.

\textbf{Confidential inference and TEEs.} Trusted execution environments and confidential accelerators protect runtime secrets, but high-throughput LLM serving often externalizes large KV state for memory efficiency. \sys{} complements these approaches by reducing what must remain secret after a cache leaves the protected runtime: copied cache tensors remain bound to the originating session basis or affine mask. GPU TEEs such as Graviton~\cite{graviton} protect runtime secrets inside an isolated execution context, whereas \sys{} targets the complementary problem of state that has already left that context.

\section{Conclusion}
\label{sec:conclusion}

We presented \sys{}, a PUF-bound KV-cache design study that separates device/session binding from full cache confidentiality. Orthogonal attention bases store K/V tensors in a device-session coordinate system while preserving attention equivalence for the legitimate device; on Qwen3-0.6B, this reduces attacker-view V-inversion from 1.000 to 0.000 and exact PII keyword leakage from 52.2\% to 0.0\%, and a migration test shows that a copied cache plus the full software stack stays uninterpretable on a device with a different PUF root, where a software-key obfuscator would not. The same orthogonal design fails against the whole basis-invariant class: norm, Gram-matrix, and singular-spectrum matching each recover protected secrets with top-1 accuracy 1.000 on Qwen3-0.6B and Llama-3.2-1B. A no-sidecar affine-mask redesign removes the class, driving every invariant attack to the $1/32$ chance level while staying fp32-exact on perplexity, HellaSwag, MMLU, and long decoding, at the cost of an in-attention de-masking step whose fused split-K Triton kernel adds $\approx 14$ percentage points of decode overhead over the orthogonal path ($-29\%$ versus $-15\%$ from plain). These results refine the core insight: the value of physical binding is non-migratability with no extractable key, and a confidential design must additionally remove the basis-invariant statistics that remain comparable after migration.
